\documentclass[11pt]{article}
\usepackage[margin=1in]{geometry}
\usepackage[T1]{fontenc}
\usepackage[utf8]{inputenc}
\usepackage{lmodern}
\usepackage[hidelinks]{hyperref}
\usepackage{array}
\usepackage{tabularx}
\usepackage{longtable}
\usepackage{booktabs}
\usepackage{enumitem}
\usepackage{titlesec}
\usepackage{xcolor}
\usepackage{fancyhdr}
\usepackage{setspace}
\usepackage{ragged2e}
\usepackage{xurl}
\setlength{\headheight}{14pt}
\pagestyle{fancy}
\fancyhf{}
\lhead{Application Security Program Handout}
\rhead{\thepage}
\renewcommand{\headrulewidth}{0.4pt}
\setlist[itemize]{leftmargin=1.25em,itemsep=0.3em,topsep=0.4em}
\setlist[enumerate]{leftmargin=1.4em,itemsep=0.35em,topsep=0.4em}
\titleformat{\section}{\large\bfseries}{\thesection.}{0.5em}{}
\titleformat{\subsection}{\normalsize\bfseries}{\thesubsection.}{0.5em}{}
\hypersetup{
  pdftitle={Building a Modern Application Security Program},
  pdfauthor={OpenAI},
  pdfsubject={Standards-based application security program handout},
  pdfcreator={LaTeX}
}
\title{\vspace{-1.2em}Building a Modern Application Security Program\\\large A standards-based handout for application and product security}
\author{}
\date{May 2026}
\emergencystretch=2em

\begin{document}
\maketitle
\thispagestyle{fancy}
\vspace{-2.0em}

\begin{abstract}
\noindent This handout summarizes a practical operating model for building an application security program using primary and authoritative sources only. The guidance is anchored in NIST's Secure Software Development Framework (SSDF), OWASP SAMM, OWASP verification and design resources, ISO/IEC 27034, and CISA secure-by-design and software supply chain guidance. It is intended for teams that need an actionable program structure without treating application security as a separate, last-minute gate.\cite{nist-ssdf,owasp-samm,iso-27034,cisa-sbd}
\end{abstract}

\section{What the program is meant to accomplish}
An application security program exists to integrate security into how software is planned, designed, built, verified, released, and maintained. NIST SP 800-218 defines the SSDF as a core set of high-level secure software development practices that can be integrated into any SDLC implementation. OWASP SAMM complements that by giving organizations a measurable maturity model for formulating and improving a software security strategy tailored to their risks. ISO/IEC 27034 frames application security as something integrated into the processes used to manage applications, whether software is built in-house, acquired, or outsourced.\cite{nist-ssdf,owasp-samm,iso-27034}

\section{Program design principles}
A defensible program should be built on the following principles:
\begin{itemize}
  \item \textbf{Security is integrated into delivery, not bolted on later.} SSDF assumes secure development practices are embedded throughout the lifecycle rather than appended to a single review stage.\cite{nist-ssdf}
  \item \textbf{Work is prioritized based on risk.} OWASP SAMM is explicitly risk-driven, and OWASP Top 10 describes itself as an awareness document and starting point rather than a complete program.\cite{owasp-samm,owasp-top10-program}
  \item \textbf{Verification criteria are explicit.} OWASP ASVS provides a basis for testing web application technical security controls and a shared set of secure development requirements.\cite{owasp-asvs}
  \item \textbf{Security ownership is distributed.} OWASP's Security Champions Guide positions champions programs as a structured way to improve security capability within delivery teams.\cite{owasp-champions}
  \item \textbf{Secure defaults matter at the product level.} CISA's secure-by-design guidance emphasizes reducing customer burden and building products that are secure to use out of the box.\cite{cisa-sbd}
\end{itemize}

\section{Core capabilities of the operating model}
\renewcommand{\arraystretch}{1.2}
\begin{longtable}{>{\RaggedRight\arraybackslash}p{0.22\textwidth} >{\RaggedRight\arraybackslash}p{0.50\textwidth} >{\RaggedRight\arraybackslash}p{0.18\textwidth}}
\toprule
\textbf{Capability} & \textbf{What the program should do} & \textbf{Primary anchors} \\
\midrule
\endfirsthead
\toprule
\textbf{Capability} & \textbf{What the program should do} & \textbf{Primary anchors} \\
\midrule
\endhead
Asset and scope definition & Identify the software, services, components, and data flows that fall within program scope so security work can be prioritized and assigned. & SSDF; ISO/IEC 27034 \\
Secure design and threat modeling & Define security requirements early, analyze architecture and trust boundaries, and identify threats before implementation and deployment. & SSDF; OWASP Threat Modeling Cheat Sheet \\
Secure implementation & Establish secure coding expectations, approved libraries and frameworks, secrets handling, and developer guidance aligned to common control requirements. & SSDF; OWASP ASVS; OWASP Top 10 \\
Verification and testing & Use automated and manual verification activities appropriate to the system, including code review, static and dynamic analysis, security testing, and requirements-based verification. & SSDF; OWASP ASVS \\
Vulnerability management & Triage findings, prioritize remediation, define fix expectations, and track closure with evidence rather than ad hoc ticket movement. & SSDF; OWASP SAMM \\
Security champions and enablement & Build team-level security capability through champions, training, reusable patterns, and collaboration models that reduce dependence on a central security bottleneck. & OWASP Security Champions Guide; OWASP SAMM \\
Supply chain and component governance & Manage dependency and component risk, generate or consume SBOM information where appropriate, and integrate supply chain controls into CI/CD and procurement workflows. & NIST supply chain guidance; CISA SBOM \\
Product security and secure defaults & Reduce exploitable default states, unsafe configurations, and customer-side hardening burden by adopting secure-by-design product decisions. & CISA Secure by Design \\
Metrics and continuous improvement & Measure whether practices are operating, whether remediation is timely, and whether program capability is maturing over time. & SSDF; OWASP SAMM \\
\bottomrule
\end{longtable}

\section{A practical phased roadmap}
\subsection{Phase 1: establish minimum program control}
\begin{enumerate}
  \item Create or validate the application and service inventory.
  \item Define minimum policy expectations for secure development, defect handling, secrets management, dependency use, and release criteria.
  \item Select a baseline verification standard for the application class, such as OWASP ASVS for web applications.
  \item Introduce core pipeline checks and a triage process for security findings.
\end{enumerate}
This phase aligns with SSDF's emphasis on preparing the organization, protecting software, and producing well-secured software, while giving the program a documented baseline for how work is supposed to happen.\cite{nist-ssdf,owasp-asvs}

\subsection{Phase 2: integrate security into engineering flow}
\begin{enumerate}
  \item Add threat modeling to design and major change workflows.
  \item Stand up a security champions model inside engineering teams.
  \item Map verification depth to application criticality and exposure.
  \item Standardize remediation, exception handling, and evidence collection.
\end{enumerate}
This phase is where security stops acting primarily as a gate and starts operating as an engineering function distributed across teams, which is consistent with both OWASP SAMM and the OWASP Security Champions guidance.\cite{owasp-samm,owasp-champions,owasp-threat-modeling}

\subsection{Phase 3: mature into product security and supply chain governance}
\begin{enumerate}
  \item Extend controls to third-party components, build provenance, and software supply chain risk management.
  \item Use secure-by-design criteria to reduce insecure default states and customer hardening burden.
  \item Expand the program for AI-related development where relevant using NIST's SSDF augmentation for generative AI and dual-use foundation models.
\end{enumerate}
This phase reflects the broader product security posture expected in modern environments, especially where cloud services, open-source dependencies, and AI-enabled systems are part of delivery.\cite{cisa-sbd,nist-218a,cisa-sbom,nist-supplychain}

\section{Metrics that are worth tracking}
Metrics should show both operational health and capability maturity. A concise set includes:
\begin{itemize}
  \item coverage of applications, repositories, or services under program scope;
  \item percentage of high-risk changes that received threat modeling;
  \item remediation time by severity and by business criticality;
  \item rate of exception use and age of unresolved exceptions;
  \item percentage of teams with an active security champion;
  \item percentage of internet-facing or high-value applications mapped to a verification baseline such as ASVS;
  \item percentage of critical products or services with software component transparency artifacts, such as SBOMs, where applicable.
\end{itemize}
OWASP SAMM is especially useful for tracking maturity over time, while SSDF is useful for checking whether core practices are actually implemented. SBOM-related metrics are supported by CISA's position that SBOMs are a key building block in software security and software supply chain risk management.\cite{owasp-samm,nist-ssdf,cisa-sbom}

\section{What not to do}
A weak program usually shows one or more of the following failure patterns:
\begin{itemize}
  \item treating the OWASP Top 10 as the entire program instead of as an awareness aid;
  \item relying only on tools without clear verification criteria, remediation ownership, or exception governance;
  \item making security a separate late-stage review that predictably conflicts with delivery timelines;
  \item focusing on application code while ignoring third-party components, build pipelines, and product defaults.
\end{itemize}
These patterns conflict with the structure implied by SSDF, SAMM, secure-by-design guidance, and current supply chain guidance.\cite{owasp-top10-program,nist-ssdf,owasp-samm,cisa-sbd,nist-supplychain}

\section{Bottom line}
A modern application security program is not a single tool stack or a one-time assessment. It is an operating model that embeds secure development practices in the SDLC, uses explicit verification criteria, distributes ownership into engineering teams, and extends beyond code-level defects into product security and software supply chain governance. The combination of NIST SSDF, OWASP SAMM, OWASP ASVS and design resources, ISO/IEC 27034, and CISA secure-by-design guidance provides a strong authoritative basis for that model.\cite{nist-ssdf,owasp-samm,owasp-asvs,iso-27034,cisa-sbd}

\section*{Authoritative sources}
\begingroup\sloppy
\begin{thebibliography}{99}
\bibitem{nist-ssdf}
National Institute of Standards and Technology (NIST), \emph{Secure Software Development Framework (SSDF) Version 1.1}, SP 800-218, 2022. Available: \url{https://csrc.nist.gov/pubs/sp/800/218/final}.

\bibitem{nist-218a}
National Institute of Standards and Technology (NIST), \emph{Secure Software Development Practices for Generative AI and Dual-Use Foundation Models}, SP 800-218A, 2024. Available: \url{https://csrc.nist.gov/pubs/sp/800/218/a/final}.

\bibitem{nist-supplychain}
National Institute of Standards and Technology (NIST), \emph{Software Security in Supply Chains}, 2022. Available: \url{https://csrc.nist.gov/pubs/other/2022/05/05/software-security-in-supply-chains/final}.

\bibitem{owasp-samm}
OWASP Foundation, \emph{OWASP SAMM}. Available: \url{https://owasp.org/www-project-samm/}.

\bibitem{owasp-champions}
OWASP Foundation, \emph{OWASP Security Champions Guide}. Available: \url{https://owasp.org/www-project-security-champions-guidebook/}.

\bibitem{owasp-threat-modeling}
OWASP Foundation, \emph{Threat Modeling Cheat Sheet}. Available: \url{https://cheatsheetseries.owasp.org/cheatsheets/Threat_Modeling_Cheat_Sheet.html}.

\bibitem{owasp-asvs}
OWASP Foundation, \emph{OWASP Application Security Verification Standard (ASVS)}. Available: \url{https://owasp.org/www-project-application-security-verification-standard/}.

\bibitem{owasp-top10-program}
OWASP Foundation, \emph{Establishing a Modern Application Security Program}, OWASP Top 10:2025. Available: \url{https://owasp.org/Top10/2025/0x03_2025-Establishing_a_Modern_Application_Security_Program/}.

\bibitem{iso-27034}
International Organization for Standardization (ISO), \emph{ISO/IEC 27034-1:2011 - Information technology - Security techniques - Application security - Part 1: Overview and concepts}. Available: \url{https://www.iso.org/standard/44378.html}.

\bibitem{cisa-sbd}
Cybersecurity and Infrastructure Security Agency (CISA), \emph{Secure by Design}. Available: \url{https://www.cisa.gov/securebydesign}.

\bibitem{cisa-sbom}
Cybersecurity and Infrastructure Security Agency (CISA), \emph{Software Bill of Materials (SBOM)}. Available: \url{https://www.cisa.gov/topics/information-communications-technology-supply-chain-security/sbom}.
\end{thebibliography}
\endgroup

\end{document}
